\documentclass[10pt,a4paper]{amsart}
\usepackage{amsmath,amssymb,amsthm,graphicx}
% Retained minimal preamble for the Stacks Project chapter sites-cohomology.tex (Cohomology on Sites).
% Verbatim excerpt of the chapter preamble: Xy-pic setup (native preamble lines 14-18) and the native
% theorem declarations with the mathematical macros (native preamble lines 154-217). Copied mechanically;
% nothing added, removed or redefined. The wrapper class and the amsmath/amssymb/amsthm/graphicx base
% packages are supplied by the assembler, so the class/externaldocument part of the native preamble is omitted.
% For commutative diagrams we use Xy-pic
\usepackage[all]{xy}

% We use 2cell for 2-commutative diagrams.
\xyoption{2cell}
% Theorem environments.
%
\theoremstyle{plain}
\newtheorem{theorem}[subsection]{Theorem}
\newtheorem{proposition}[subsection]{Proposition}
\newtheorem{lemma}[subsection]{Lemma}

\theoremstyle{definition}
\newtheorem{definition}[subsection]{Definition}
\newtheorem{example}[subsection]{Example}
\newtheorem{exercise}[subsection]{Exercise}
\newtheorem{situation}[subsection]{Situation}

\theoremstyle{remark}
\newtheorem{remark}[subsection]{Remark}
\newtheorem{remarks}[subsection]{Remarks}

\numberwithin{equation}{subsection}

% Macros
%
\def\lim{\mathop{\mathrm{lim}}\nolimits}
\def\colim{\mathop{\mathrm{colim}}\nolimits}
\def\Spec{\mathop{\mathrm{Spec}}}
\def\Hom{\mathop{\mathrm{Hom}}\nolimits}
\def\Ext{\mathop{\mathrm{Ext}}\nolimits}
\def\SheafHom{\mathop{\mathcal{H}\!\mathit{om}}\nolimits}
\def\SheafExt{\mathop{\mathcal{E}\!\mathit{xt}}\nolimits}
\def\Sch{\mathit{Sch}}
\def\Mor{\mathop{\mathrm{Mor}}\nolimits}
\def\Ob{\mathop{\mathrm{Ob}}\nolimits}
\def\Sh{\mathop{\mathit{Sh}}\nolimits}
\def\NL{\mathop{N\!L}\nolimits}
\def\CH{\mathop{\mathrm{CH}}\nolimits}
\def\proetale{{pro\text{-}\acute{e}tale}}
\def\etale{{\acute{e}tale}}
\def\QCoh{\mathit{QCoh}}
\def\Ker{\mathop{\mathrm{Ker}}}
\def\Im{\mathop{\mathrm{Im}}}
\def\Coker{\mathop{\mathrm{Coker}}}
\def\Coim{\mathop{\mathrm{Coim}}}

% Boxtimes
%
\DeclareMathSymbol{\boxtimes}{\mathbin}{AMSa}{"02}

%
% Macros for moduli stacks/spaces
%
\def\QCohstack{\mathcal{QC}\!\mathit{oh}}
\def\Cohstack{\mathcal{C}\!\mathit{oh}}
\def\Spacesstack{\mathcal{S}\!\mathit{paces}}
\def\Quotfunctor{\mathrm{Quot}}
\def\Hilbfunctor{\mathrm{Hilb}}
\def\Curvesstack{\mathcal{C}\!\mathit{urves}}
\def\Polarizedstack{\mathcal{P}\!\mathit{olarized}}
\def\Complexesstack{\mathcal{C}\!\mathit{omplexes}}
% \Pic is the operator that assigns to X its picard group, usage \Pic(X)
% \Picardstack_{X/B} denotes the Picard stack of X over B
% \Picardfunctor_{X/B} denotes the Picard functor of X over B
\def\Pic{\mathop{\mathrm{Pic}}\nolimits}
\def\Picardstack{\mathcal{P}\!\mathit{ic}}
\def\Picardfunctor{\mathrm{Pic}}
\def\Deformationcategory{\mathcal{D}\!\mathit{ef}}
\usepackage{cleveref}
\crefname{originallemma}{original lemma}{original lemmas}
\Crefname{originallemma}{Original lemma}{Original lemmas}
\makeatletter
\newlabel{lemma-restrict-K-injective-to-open}{{08FI}{source}{}{sourcepointer-1}{}}
\newlabel{lemma-restrict-K-injective-to-open@cref}{{[originallemma][0][]08FI}{[1][0][]source}{}{}{}}
\makeatother
\crefname{originallemma}{original lemma}{original lemmas}
\Crefname{originallemma}{Original lemma}{Original lemmas}
\makeatletter
\newlabel{sites-modules-lemma-extension-by-zero}{{03DI}{source}{}{sourcepointer-2}{}}
\newlabel{sites-modules-lemma-extension-by-zero@cref}{{[originallemma][0][]03DI}{[1][0][]source}{}{}{}}
\makeatother
\crefname{originallemma}{original lemma}{original lemmas}
\Crefname{originallemma}{Original lemma}{Original lemmas}
\makeatletter
\newlabel{derived-lemma-K-injective}{{070I}{source}{}{sourcepointer-3}{}}
\newlabel{derived-lemma-K-injective@cref}{{[originallemma][0][]070I}{[1][0][]source}{}{}{}}
\makeatother
\title{Extension by zero is left adjoint to restriction on the derived category}
\author{The Stacks Project Authors}
\date{stacks-project:sites-cohomology.tex:local-snapshot-20261011; GNU FDL 1.2 or later}
\begin{document}
\maketitle
\noindent\textit{Editorial note.} All delivered mathematical text is the author's own, copied from the pinned Stacks Project chapter file sites-cohomology.tex; no statement or proof step was rewritten, reordered or removed.\par\medskip
\noindent\textit{Editorial note.} The proof is the author's complete argument; the abelian adjunction and the passage to the derived categories are both inside the counted span.\par\medskip
\noindent\textit{Editorial note.} Every labelled reference occurring inside the counted spans is carried in assembly.external\_references as a non-load-bearing navigation pointer (same-chapter ones at their definition site, cross-chapter ones at the invocation inside the counted span), so no load-bearing mathematics is replaced by a stub. No bibliography citation occurs inside the counted spans.\par\medskip
\noindent\textit{Editorial note.} The checked chapter list for same-name collisions is the set of stacks chapters staged in this repository: chow, cohomology, divisors, etale-cohomology, flat, more-algebra, more-morphisms, morphisms, spaces-more-morphisms, varieties together with this card's chapter sites-cohomology.\par\medskip
\noindent\textit{Editorial note.} The preamble is the chapter's own retained minimal preamble (sources/stacks-sites-cohomology/retained-preamble.tex), a verbatim excerpt of the native preamble: the Xy-pic setup and the native theorem declarations with the mathematical macros.\par\medskip
\noindent\textit{Allowed standard prerequisites (editorial).}\begin{itemize}
\item Extension by zero j\_! and restriction along the localization morphism. Same chapter, the functors j\_! and the restriction on modules over a ringed site and their adjunction on the abelian level. Supplies the functors whose derived adjunction is asserted.
\item Derived functors of adjoint pairs and the computation of the derived adjunction. Same chapter together with the chapter Derived Categories, the passage from an adjunction of abelian categories, one side exact, to an adjunction of the derived categories. Supplies the argument upgrading the adjunction to the derived categories.
\end{itemize}\medskip
\ifdefined\hypertarget\hypertarget{sourcepointer-1}{}\fi
\noindent\textit{Source reference.} Original lemma 08FI: sources/stacks-sites-cohomology/pinned-original-source.tex, source lines 4314--4314. Navigation pointer invoked at proof line 4564; the counted span contains the pointer only and the cited result is not reproduced inside it, so the pointer is registered as a non-load-bearing navigation stub in the established form for Stacks material. No bibliography citation is involved.\par\medskip
\ifdefined\hypertarget\hypertarget{sourcepointer-2}{}\fi
\noindent\textit{Source reference.} Original lemma 03DI: sources/stacks-sites-cohomology/pinned-original-source.tex, source lines 4572--4572. Navigation pointer invoked at proof line 4572; the counted span contains the pointer only and the cited result is not reproduced inside it, so the pointer is registered as a non-load-bearing navigation stub in the established form for Stacks material. No bibliography citation is involved.\par\medskip
\ifdefined\hypertarget\hypertarget{sourcepointer-3}{}\fi
\noindent\textit{Source reference.} Original lemma 070I: sources/stacks-sites-cohomology/pinned-original-source.tex, source lines 4574--4574. Navigation pointer invoked at proof line 4574; the counted span contains the pointer only and the cited result is not reproduced inside it, so the pointer is registered as a non-load-bearing navigation stub in the established form for Stacks material. No bibliography citation is involved.\par\medskip
\begin{lemma}
\label{lemma-adjoint-lower-shriek-restrict}
Let $(\mathcal{C}, \mathcal{O})$ be a ringed site. Let $U$ be an
object of $\mathcal{C}$. Denote
$j : (\Sh(\mathcal{C}/U), \mathcal{O}_U) \to (\Sh(\mathcal{C}), \mathcal{O})$
the corresponding localization morphism. The restriction functor
$D(\mathcal{O}) \to D(\mathcal{O}_U)$ is a right adjoint to
extension by zero $j_! : D(\mathcal{O}_U) \to D(\mathcal{O})$.
\end{lemma}

\begin{proof}
We have to show that
$$
\Hom_{D(\mathcal{O})}(j_!E, F) = \Hom_{D(\mathcal{O}_U)}(E, F|_U)
$$
Choose a complex $\mathcal{E}^\bullet$ of $\mathcal{O}_U$-modules
representing $E$ and choose
a K-injective complex $\mathcal{I}^\bullet$ representing $F$.
By Lemma \ref{lemma-restrict-K-injective-to-open} the complex
$\mathcal{I}^\bullet|_U$ is K-injective as well. Hence we see that
the formula above becomes
$$
\Hom_{D(\mathcal{O})}(j_!\mathcal{E}^\bullet, \mathcal{I}^\bullet) =
\Hom_{D(\mathcal{O}_U)}(\mathcal{E}^\bullet, \mathcal{I}^\bullet|_U)
$$
which holds as $|_U$ and $j_!$ are adjoint functors
(Modules on Sites, Lemma \ref{sites-modules-lemma-extension-by-zero})
and
Derived Categories, Lemma \ref{derived-lemma-K-injective}.
\end{proof}

\end{document}
